\chapter{Two-Dimensional Supersymmetric Models}
\label{ch:susy-two-dimensional-lg_sigma_glsm}

\ConventionsMixed

This chapter begins the lower-dimensional supersymmetric examples.  The main
objects are \(2D\) \(\mathcal N=(2,2)\) Landau--Ginzburg models, sigma
models, and gauged linear sigma models.  Each is treated as a QFT datum with
declared fields, symmetries, operator coordinates, and infrared questions.

\section{The N=(2,2) Algebra}

In Lorentzian signature let \(P_\pm=P_0\pm P_1\).  The supercharges
\[
  Q_\pm,\quad \bar Q_\pm
\]
obey
\[
  \{Q_+,\bar Q_+\}=2P_+,\qquad
  \{Q_-,\bar Q_-\}=2P_-,
\]
with all other anticommutators vanishing in the absence of central charges.
An \(R\)-symmetry assignment consists of vector and axial charges
\[
  F_V,\qquad F_A,
\]
whose action on supercharges is part of the model datum.  A BPS state with
central charge \(Z\) obeys an inequality obtained by positivity of
\((Q+\ee^{\ii\alpha}\bar Q)^\dagger(Q+\ee^{\ii\alpha}\bar Q)\); the saturated
states are kernels of a corresponding linear combination of supercharges.

\section{Landau--Ginzburg Datum}

A Landau--Ginzburg model with chiral superfields \(\Phi^i\) is specified by
a target \(X=\C^n\) or a complex manifold chart, a Kahler potential \(K\), and
a holomorphic superpotential
\[
  W:X\longrightarrow\C .
\]
In components, after eliminating auxiliary fields, the bosonic potential is
\[
  V(\phi)=g^{i\bar j}\partial_iW(\phi)\,
  \overline{\partial_jW(\phi)}.
\]
Supersymmetric classical vacua are critical points
\[
  \partial_iW(\phi)=0.
\]
When the critical points are isolated and nondegenerate, the chiral ring is
\[
  \mathcal R_W
  =
  \C[\phi^1,\ldots,\phi^n]\big/
  \bigl(\partial_1W,\ldots,\partial_nW\bigr).
\]
The multiplication is ordinary polynomial multiplication followed by quotient.
The residue pairing is
\[
  \langle f,g\rangle_{\mathrm{res}}
  =
  \sum_{\partial W=0}
  \frac{f(\phi)g(\phi)}
  {\det(\partial_i\partial_jW)(\phi)} ,
\]
provided all critical points are nondegenerate.  If degeneracies occur, the
Jacobian algebra remains the algebraic chiral ring, while the analytic
definition of the pairing requires a contour or local-duality prescription.

\section{B-Type Topological Sector}

The B-twist uses a scalar supercharge \(Q_B\).  The \(Q_B\)-closed local
operators in a Landau--Ginzburg model are represented by holomorphic
functions modulo \(Q_B\)-exact terms.  The auxiliary-field equations imply
\[
  \partial_iW(\phi)=Q_B(\cdots)_i ,
\]
inside \(Q_B\)-cohomology, hence the local operator algebra is \(\mathcal R_W\).
On a compact Riemann surface the finite-dimensional localization formula
takes the form
\[
  \langle f_1\cdots f_m\rangle_{g}
  =
  \sum_{\partial W=0}
  \frac{f_1(\phi)\cdots f_m(\phi)}
  {\det(\partial_i\partial_jW)(\phi)^{1-g}},
\]
after a normalization of zero modes and determinants is fixed.  This formula
is a theorem of the finite-dimensional localized model; the QFT statement
requires a regulator preserving the chosen \(Q_B\) and a proof that no
boundary terms are produced in the localization deformation.

\section{Sigma-Model Datum}

A \(2D\) \(\mathcal N=(2,2)\) nonlinear sigma model is specified by a Kahler
target \(X\), metric \(g_{i\bar j}\), \(B\)-field \(B\), and possibly a
superpotential \(W\).  The bosonic Euclidean action is
\[
  S_{\mathrm{bos}}
  =
  \frac1{4\pi\alpha'}
  \int_\Sigma
  \dd^2x\,
  g_{i\bar j}(\phi)\partial_\mu\phi^i\partial_\mu\bar\phi^{\bar j}
  +
  \frac{\ii}{2\pi\alpha'}
  \int_\Sigma \phi^*B .
\]
The large-volume perturbative beta tensor begins with the Ricci form:
\[
  \mu\frac{\dd g_{i\bar j}}{\dd\mu}
  =
  \frac1{2\pi}R_{i\bar j}+\cdots .
\]
Thus Ricci-flatness is the first condition for a conformal infrared candidate
in a large-volume chart.  The full QFT claim includes instanton corrections,
operator renormalization, and a continuum limit; these are additional data,
not consequences of the local tensor equation alone.

\section{Gauged Linear Sigma-Model Datum}

A GLSM is specified by a compact gauge group \(G\), chiral multiplets in a
complex representation \(V\), a superpotential \(W\in\mathcal O(V)^G\), and
Fayet--Iliopoulos-theta parameters for abelian factors.  For
\[
  G=U(1),\qquad \Phi_i\text{ of charge }Q_i,
\]
the classical D-term equation is
\[
  \sum_i Q_i|\phi_i|^2=r .
\]
The quotient by \(U(1)\) gives a Kahler quotient phase when \(r\) lies in a
chamber with stable points.  The axial \(R\)-symmetry anomaly coefficient is
\[
  \sum_i Q_i .
\]
When this coefficient vanishes, the FI parameter is not shifted by the axial
rotation at one loop, and the GLSM is a candidate ultraviolet presentation of
a Calabi--Yau sigma model.  The precise infrared statement is a comparison
between two QFTs: a gauge-theory continuum limit and a sigma-model continuum
limit in the chamber under consideration.

\section{Infrared Comparison Problem}

For each two-dimensional supersymmetric model one must distinguish:
\[
\begin{gathered}
  \text{UV regulator datum},\qquad
  \text{supersymmetric Wilsonian flow},\\
  \text{protected operator algebra},\qquad
  \text{full infrared QFT}.
\end{gathered}
\]
Protected rings and indices can match before the full Hilbert space,
correlators, and stress tensor are identified.  A complete infrared
equivalence therefore requires a map of local algebras, states, stress
tensors, supercharges, and extended operators, with convergence or scaling
limits stated.

\begin{openproblem}[Two-dimensional supersymmetric RG construction]
\label{op:two-dimensional-susy-rg-construction}
Construct \(2D\) \(\mathcal N=(2,2)\) GLSM flows to Landau--Ginzburg and
Calabi--Yau sigma-model fixed points with a regulator, supersymmetric
Wilsonian maps, protected-ring comparison, stress-tensor convergence, and
Hilbert-space reconstruction in one declared framework.
\end{openproblem}
